\documentclass[12pt]{article}
\usepackage[utf8]{inputenc}
\usepackage[T1]{fontenc}
\usepackage{hyperref}
\title{Yug Language Resources for K-12 Curriculum Development}
\author{Compiled by Claude Sonnet 5.5 (claude-sonnet-5-5) for Andr\'as Kornai}
\date{}
\begin{document}
\maketitle

\section{Introduction}
Yug (Yugh, formerly called Sym Ket or Southern Ket; ISO 639-3 \texttt{yug}, Glottocode yugh1239) is an extinct Yeniseian language of Siberia, closely related to Ket. Search results report that it went extinct by about 1972. Heinrich Werner documented it with the last speakers between 1961 and 1971. There are no native speakers, so a curriculum would be a heritage and revitalization effort built on archival and scholarly documentation. I found no school programmes, primers or audio for Yug in my searches.

\section{Primary Educational Resources}
I found no primers, textbooks or classroom materials for Yug.

\section{Linguistic and Reference Materials}
\begin{itemize}
\item Wikipedia, ``Yugh language'': overview of the language and its history. \url{https://en.wikipedia.org/wiki/Yugh_language}
\item Wikipedia, ``Heinrich Werner (linguist)'': background on the main documenter of Yug. \url{https://en.wikipedia.org/wiki/Heinrich_Werner_(linguist)}
\item Glottolog entry for Yug, with bibliographic references. The search summary mentions Werner's \emph{Das Jugische (Sym-Ketische)} (1997) and the \emph{Dictionary of the Yugh Language} (Lincom, 2012, ed.\ Edward J.\ Vajda). \url{https://glottolog.org/resource/languoid/iso/yug}
\item Edward Vajda, publication listing on bepress, which came up in the search for Yug materials. \url{https://works.bepress.com/edward_vajda/74}
\item Paper from the Heidelberg ICHL 26 conference, returned by the search. I did not confirm its exact topic. \url{https://www.uni-heidelberg.de/md/slav/forschung/tagungen/ichl26/ichl26_paper_32.pdf}
\item LINGUIST List issue 34.65, returned by the search. I did not confirm its exact contents. \url{https://dev.linguistlist.org/issues/34/65/}
\end{itemize}

\section{Religious and Community Resources}
I found no religious or community resources for Yug in my searches.

\section{Digital and Audio Resources}
I found no audio or digital teaching resources for Yug. The search did not surface a digitized version of Werner's dictionary.

\section{Curriculum Development Context}
Yug has no living community and no known school programme. Werner's dictionary and grammar are the core documentation. Any curriculum would have to be built from these scholarly sources, ideally with Yeniseian specialists such as Edward Vajda. A related option is the living Ket language, though the sources found here do not describe Ket teaching materials.

\end{document}